\clearpage
\VerdierPhysicalPageAnchor{0105}%
\begingroup
\small
\setlength{\abovedisplayskip}{3pt plus 2pt minus 2pt}%
\setlength{\belowdisplayskip}{3pt plus 2pt minus 2pt}%
\setlength{\jot}{2pt}%

\VerdierSemanticAnchor{ch1-3-4-3-cone-of-hf}
On the other hand, one has (definition of the cone functor):
\[
  c(H(f))^l=H(Y)^l\oplus(H(X)[1])^l,
\]

\VerdierSemanticAnchor{ch1-3-4-3-4}
\begin{equation}
  d_{c(H(f))}^{,l}=
  \begin{pmatrix}
    d_{H(Y)}^{,l} & H(f)^{l+1}\\
    0 & d_{H(X)[1]}^{,l}
  \end{pmatrix},
  \tag{3.4.3.4}
\end{equation}
where the morphism $H(f)$ is defined by (2.4.2.6), (2.4.2.3), and
(2.2.2.6):

\VerdierSemanticAnchor{ch1-3-4-3-5}
\begin{equation}
  H(f)^{l+1}=
  \prod_{\substack{\sigma=\sum_i l_i e_i\\ \sum_i l_i=l}}
  F(X_1^{l_1},\ldots,f^{l_j+1},\ldots,X_n^{l_n}).
  \tag{3.4.3.5}
\end{equation}

\VerdierSemanticAnchor{ch1-3-4-3-j-map}
Now set
\[
  J^l\colon H(Y)^l\oplus H(X[1])^l
  \longrightarrow H(Y)^l\oplus(H(X)[1])^l,
\]

\VerdierSemanticAnchor{ch1-3-4-3-6}
\begin{equation}
  J^l=
  \begin{pmatrix}
    \operatorname{id}_{H(Y)^l} & 0\\
    0 & \Psi^l
  \end{pmatrix},
  \qquad
  \Psi^l=
  \prod_{\substack{\sigma=\sum_i l_i e_i\\ \sum_i l_i=l}}
  \varepsilon(\sigma,j)
  \operatorname{id}_{F(X_1^{l_1},\ldots,X^{l_j+1},\ldots,X_n^{l_n})}.
  \tag{3.4.3.6}
\end{equation}

\VerdierSemanticAnchor{ch1-3-4-3-7}
To verify that the $J^l$ are the components of a morphism of complexes, it
suffices to verify that
\begin{equation}
  \Phi^l=H(f)^{l+1}\Psi^l,
  \qquad
  \Psi^{l+1}d_{H(X[1])}^{l}=d_{H(X)[1]}^{l}\Psi^l.
  \tag{3.4.3.7}
\end{equation}

\VerdierSemanticAnchor{ch1-3-4-3-8}
To verify that diagram (3.4.3.1) is commutative, it suffices to verify that
\begin{equation}
  \Psi^l=m^l.
  \tag{3.4.3.8}
\end{equation}
\endgroup
